%==============================================================================
%  5. Conclusion
%==============================================================================

\section{Conclusion}
\label{sec:conclusion}

This paper presented an LLM-enhanced differential evolution framework for
multi-UAV cooperative path planning in which a pretrained language model acts as
a plug-in controller for the three search mechanisms of DMDE -- the crossover
rate, the scale factor and the extinction rate -- which are otherwise coupled by
hand-designed equations, leaving the evolutionary operators unchanged. The
framework was evaluated through an ablation study on a balanced and a
scarce-resource scenario.
% 待补：一句话概括最重要的结论（务必如实：离线固定 CR=0.3 在 9/9 组对照中胜出）
\todo{summarise the single most important conclusion in one sentence.}

The study also exposes a limitation that we consider practically important:
decoupling the three mechanisms does not by itself make the language model a
better controller. On this static assignment task an offline constant
$\CR = 0.3$ outperforms every online variant. We attribute this to three
factors. Operator or strategy selection is a semantic classification task that a
pretrained model can recall from its training distribution, whereas parameter
control is a numeric calibration task for which no such prior exists. The action
granularity is also unfavourable: adjacent \CR{} levels differ in effect by
about $-0.248$~pp while between-run noise has a standard deviation of
$3.513$~pp, a signal-to-noise ratio of roughly $1/14$. Finally, the available
headroom is small and the task is static, so there is no time-varying structure
for an online policy to exploit. Future work will therefore move the action
space from numeric parameters to a semantic strategy tier, letting the model
choose among qualitatively distinct search modes while a bandit performs the
numeric credit assignment. % 待补：其它未来工作方向（多目标扩展、动态场景）
\todo{add further directions, e.g. multi-objective extensions and dynamic
      scenarios with time-varying cost matrices.}
